\documentclass[12pt,a4paper]{article}
\usepackage[T2A]{fontenc}
\usepackage[utf8]{inputenc}
\usepackage[russian]{babel}
\usepackage{geometry}
\usepackage{amsmath,amssymb,mathtools}
\usepackage{array,booktabs,longtable}
\usepackage[dvipsnames]{xcolor}
\usepackage{tikz}
\usepackage{hyperref}
\geometry{left=20mm,right=20mm,top=22mm,bottom=22mm}
\hypersetup{colorlinks=true,linkcolor=blue!55!black}
\setlength{\parindent}{0pt}
\setlength{\parskip}{6pt}
\setlength{\emergencystretch}{2em}
\renewcommand{\arraystretch}{1.2}
\newcommand{\rub}{\,\text{руб.}}
\newcommand{\Err}{\textcolor{red}{\textbf{Ошибка:}}}
\newcommand{\Idea}{\textcolor{blue!60!black}{\textbf{Ключевая идея:}}}
\newcommand{\Result}{\textcolor{teal!60!black}{\textbf{Итог:}}}
\begin{document}

% Титульный лист
\thispagestyle{empty}
\vspace*{0.28\textheight}
\begin{center}
{\Huge\bfseries Ревью домашней работы\par}
\vspace{20mm}
{\large Автор: Лёвин Артём Александрович\par}
\vspace{5mm}
{\large 23 октября 2026 года\par}
\end{center}
\vfill
\begin{center}
\begin{tikzpicture}
\draw[blue!40!black,line width=0.6pt]
  (0,0) .. controls (2,0.7) and (4,-0.7) .. (6,0)
        .. controls (8,0.7) and (10,-0.7) .. (12,0);
\end{tikzpicture}
\end{center}
\newpage

\section*{Сводная таблица ошибок}
{\small
\setlength{\tabcolsep}{4pt}
\begin{longtable}{@{}p{0.9cm}p{2.5cm}p{4.65cm}p{3.0cm}p{3.5cm}@{}}
\toprule
\textbf{№} & \textbf{Тема} & \textbf{Суть ошибки} & \textbf{Ваш ответ} & \textbf{Правильный ответ} \\
\midrule
\endhead
\hyperref[task:6]{6} & Кредит с погашением долга по долям & Общая выплата ошибочно приравнена к переплате. & Ставка не найдена. & 5\% годовых. \\
\midrule
\hyperref[task:8]{8} & Равные годовые платежи & В уравнении пропущен платёж второго года; расчёт не завершён. & Не указан. & Платёж 43\,200 руб.; переплата 38\,600 руб. (42,42\%). \\
\midrule
\hyperref[task:10]{10} & Сравнение способов погашения & Найдена формула равного платежа, но не вычислены обе общие выплаты и их разность. & Вывод не указан. & Выгоднее равные доли долга на 2\,417,58 руб. \\
\bottomrule
\end{longtable}
}

\newpage
\thispagestyle{empty}
\vspace*{\fill}
\begin{center}
{\LARGE\bfseries Разбор ошибок}
\end{center}
\vspace*{\fill}
\newpage

\phantomsection\label{task:6}
\section*{Задание № 6}
\textbf{Текст задания.} При дифференцированном погашении кредита за пять лет общая переплата составила 15\% от первоначальной суммы. Найдите постоянную годовую процентную ставку.

\subsection*{Ваше решение}
Вы верно обозначили первоначальную сумму через $S$, число лет через $N=5$ и записали, что переплата составляет $0{,}15S$. Затем использовали формулу общей выплаты при дифференцированном погашении:
\[
T=S+\frac{Sp(N+1)}{2}.
\]
Однако далее эта общая выплата была приравнена к $0{,}15S$:
\[
S+\frac{Sp(N+1)}{2}=0{,}15S.
\]
\textbf{Ваш ответ:} годовая ставка не найдена.

\subsection*{Ошибка и её точное место}
\Err{} перепутаны \emph{переплата} и \emph{общая сумма выплат}.

\textbf{Последний правильный шаг:}
\[
\Pi=0{,}15S,\qquad T=S+\frac{Sp(N+1)}{2},
\]
где $\Pi$ означает переплату, а $T$ --- всё, что возвращено банку.

\textbf{Первый неверный шаг:}
\[
\underbrace{S+\frac{Sp(N+1)}{2}}_{T}=\underbrace{0{,}15S}_{\Pi}.
\]

\subsection*{Почему это неверно}
Первоначальная сумма долга $S$ не является переплатой: эти деньги банк выдал заёмщику, и их необходимо вернуть независимо от процентов. Переплата --- только проценты сверх возвращённого долга. Поэтому
\[
T=S+\Pi,\qquad \Pi=T-S.
\]
При положительной ставке общая выплата больше первоначального долга, а равенство $T=0{,}15S$ означало бы возврат лишь 15\% полученных денег. Это противоречит условию погашения кредита.

\Idea{} В формуле общей выплаты сначала вычтите первоначальный долг. Полученное выражение для процентов нужно приравнять к 15\% первоначальной суммы.

\subsection*{Исправление от последнего верного шага}
\[
\begin{aligned}
\Pi&=T-S
=\left(S+\frac{Sp(N+1)}{2}\right)-S\\
&=\frac{Sp(N+1)}{2}.
\end{aligned}
\]
По условию $\Pi=0{,}15S$, поэтому
\[
\frac{Sp(N+1)}{2}=0{,}15S.
\]

\subsection*{Полное правильное решение}
При дифференцированном погашении основную сумму долга возвращают равными долями. За пять лет ежегодно погашается одна пятая исходного долга. Перед начислением процентов остатки равны
\[
S,\quad \frac45S,\quad \frac35S,\quad \frac25S,\quad \frac15S.
\]
Обозначим годовую ставку \emph{в долях единицы} через $p$. Проценты за пять лет составят
\[
\begin{aligned}
\Pi
&=pS+\frac45pS+\frac35pS+\frac25pS+\frac15pS\\
&=pS\left(1+\frac45+\frac35+\frac25+\frac15\right)\\
&=3pS.
\end{aligned}
\]
Та же величина получается по формуле
\[
\Pi=\frac{Sp(5+1)}2=3Sp.
\]
Из условия $3Sp=0{,}15S$. Поскольку сумма кредита положительна, делим обе части на $S$:
\[
3p=0{,}15,\qquad p=0{,}05.
\]
Переводим долю единицы в проценты: $0{,}05=5\%$.

\subsection*{Проверка}
При ставке 5\% переплата равна $3\cdot0{,}05S=0{,}15S$, то есть ровно 15\% исходной суммы. Общая выплата равна $1{,}15S$, что согласуется с возвратом долга и процентов.

\Result{} постоянная годовая процентная ставка равна \textbf{5\%}.

\subsection*{Задача на закрепление}
Кредит погашается дифференцированными платежами в течение четырёх лет. Общая переплата за весь срок составила 20\% первоначальной суммы кредита. Найдите постоянную годовую процентную ставку.

\newpage
\phantomsection\label{task:8}
\section*{Задание № 8}
\textbf{Текст задания.} Кредит 91\,000 руб. взят на три года под 20\% годовых и погашается равными годовыми платежами. Найдите размер одного платежа, переплату в рублях и процент переплаты. Процент округлите до сотых.

\subsection*{Ваше решение}
Вы обозначили первоначальную сумму через $S=91\,000$, число лет через $N=3$, ставку через $p=0{,}2$, а коэффициент ежегодного увеличения долга через $k=1{,}2$. Для равного платежа $x$ Вы записали:
\[
Sk^3-xk^2-x=0,
\]
а затем получили выражение
\[
x=\frac{Sk^3}{k^2+1}
=\frac{157\,248}{2{,}44}.
\]
Вычисление платежа и расчёт переплаты не завершены.

\textbf{Ваш ответ:} не указан.

\subsection*{Ошибка и её точное место}
\Err{} не учтён платёж второго года при приведении всех сумм к концу третьего года.

\textbf{Последний правильный шаг:} определить $S=91\,000$, $k=1{,}2$, $N=3$ и обозначить одинаковые платежи через $x$.

\textbf{Первый неверный шаг:}
\[
Sk^3-xk^2-x=0.
\]
Между слагаемыми $-xk^2$ и $-x$ должно стоять ещё $-xk$.

\subsection*{Почему это неверно}
Первый платёж уменьшает долг в конце первого года. Без него задолженность за следующие два года увеличилась бы дважды, поэтому к концу третьего года влияние первого платежа равно $xk^2$.

Второй платёж вносится в конце второго года, а затем проходит ещё один год начисления процентов. Его влияние на долг к концу третьего года равно $xk$.

Третий платёж вносится в конце третьего года, поэтому его влияние равно просто $x$. Получается \emph{три} вычитаемых слагаемых, а не два:
\[
Sk^3-xk^2-xk-x=0.
\]
Если пропустить средний платёж, он как будто вообще не уменьшает кредит, поэтому размер рассчитанного платежа получится завышенным.

\Idea{} После каждого года сначала начисляйте проценты на остаток долга, затем вычитайте очередной одинаковый платёж. Это позволяет не пропустить ни один платёж.

\subsection*{Исправление от последнего верного шага}
\[
\begin{aligned}
Sk^3-xk^2-xk-x&=0,\\
Sk^3&=x(k^2+k+1),\\
x&=\frac{Sk^3}{k^2+k+1}.
\end{aligned}
\]
Здесь знаменатель равен $1{,}44+1{,}2+1=3{,}64$, а не $2{,}44$.

\subsection*{Полное правильное решение}
\textbf{Шаг 1.} Обозначим равный ежегодный платёж через $x$. Коэффициент роста задолженности за год равен $1+0{,}20=1{,}2$.

\textbf{Шаг 2.} К концу первого года после начисления процентов и платежа останется
\[
B_1=91\,000\cdot1{,}2-x=109\,200-x.
\]
\textbf{Шаг 3.} К концу второго года останется
\[
\begin{aligned}
B_2&=1{,}2B_1-x\\
&=1{,}2(109\,200-x)-x\\
&=131\,040-2{,}2x.
\end{aligned}
\]
\textbf{Шаг 4.} После третьего платежа долг должен стать равен нулю:
\[
\begin{aligned}
0=B_3&=1{,}2B_2-x\\
&=1{,}2(131\,040-2{,}2x)-x\\
&=157\,248-3{,}64x.
\end{aligned}
\]
Отсюда
\[
x=\frac{157\,248}{3{,}64}=43\,200\rub.
\]
\textbf{Шаг 5.} За три года выплачено
\[
T=3\cdot43\,200=129\,600\rub.
\]
Переплата равна общей выплате за вычетом выданной суммы:
\[
\Pi=129\,600-91\,000=38\,600\rub.
\]
Процент переплаты считают \emph{от первоначального кредита}:
\[
q=\frac{38\,600}{91\,000}\cdot100\%
\approx42{,}42\%.
\]

\subsection*{Проверка}
После первого платежа остаток равен $109\,200-43\,200=66\,000$ руб. После второго --- $66\,000\cdot1{,}2-43\,200=36\,000$ руб. После третьего --- $36\,000\cdot1{,}2-43\,200=0$. Долг погашен точно.

\Result{} ежегодный платёж --- \textbf{43\,200 руб.}; переплата --- \textbf{38\,600 руб.}, или \textbf{42,42\%} от суммы кредита.

\subsection*{Задача на закрепление}
Кредит 66\,200 руб. взят на три года под 10\% годовых и погашается тремя равными ежегодными платежами в конце каждого года. Найдите размер одного платежа, переплату в рублях и процент переплаты. Процент округлите до сотых.

\newpage
\phantomsection\label{task:10}
\section*{Задание № 10}
\textbf{Текст задания.} Человек выбирает способ погашения кредита 100\,000 руб. на три года под 20\% годовых: равными платежами или равными частями основного долга. Какой способ даёт меньшую общую выплату и на сколько рублей?

\subsection*{Ваше решение}
Вы верно обозначили $S=100\,000$, $N=3$, $p=0{,}2$, $k=1{,}2$. Для равных платежей $x$ записано корректное уравнение
\[
Sk^3-xk^2-xk-x=0
\]
и получена верная формула
\[
x=\frac{Sk^3}{k^2+k+1}
=\frac{100\,000\cdot1{,}728}{3{,}64}.
\]
Также отмечено, что общая сумма при равных платежах равна $3x$. Но численный результат, сумма при дифференцированном погашении и сравнение не приведены.

\textbf{Ваш ответ:} вывод о более выгодном способе и разности выплат не указан.

\subsection*{Ошибка и её точное место}
\Err{} решение не завершено: для выбора способа погашения требуется вычислить \emph{обе} общие суммы и их разность.

\textbf{Последний правильный шаг:}
\[
x=\frac{100\,000\cdot1{,}728}{3{,}64},\qquad T_{\text{равн.}}=3x.
\]
\textbf{Первый неполный шаг:} после нахождения выражения для $x$ расчёт прекращён, а сумма выплат при погашении равными частями основного долга не найдена.

\subsection*{Почему этого недостаточно}
Вопрос требует сравнения двух вариантов. Знание суммы платежа только в первом варианте ещё не даёт ответа, какой из двух вариантов дешевле. Кроме того, равные части основного долга --- не то же самое, что равные полные платежи: в первом случае основная часть платежа постоянна, а проценты уменьшаются вместе с остатком долга.

\Idea{} Для каждого способа отдельно найдите общую сумму, включая возвращённый долг и начисленные проценты. Затем вычтите меньшую сумму из большей.

\subsection*{Исправление от последнего верного шага}
Из уже правильно полученной формулы
\[
x=\frac{172\,800}{3{,}64}\approx47\,472{,}52747\rub.
\]
Получаем
\[
T_{\text{равн.}}=3x\approx142\,417{,}58\rub.
\]
Остаётся независимо рассчитать вторую схему.

\subsection*{Полное правильное решение}
\textbf{Способ 1. Равные полные платежи.} Ежегодно в конце года платят одну и ту же сумму. Из правильного уравнения Вашего решения получаем
\[
\begin{aligned}
x&=\frac{100\,000\cdot1{,}2^3}{1{,}2^2+1{,}2+1}\\
&=\frac{172\,800}{3{,}64}
\approx47\,472{,}52747\rub.
\end{aligned}
\]
Общая сумма при вычислении без промежуточного округления:
\[
T_1=3x\approx142\,417{,}58\rub.
\]

\textbf{Способ 2. Равные части основного долга.} Каждый год возвращается треть исходной суммы:
\[
a=\frac{100\,000}{3}\rub.
\]
В первом году проценты начисляются на все $100\,000$ руб.:
\[
I_1=100\,000\cdot0{,}2=20\,000\rub.
\]
Во втором году осталось две трети основного долга, поэтому
\[
I_2=\frac{2\cdot100\,000}{3}\cdot0{,}2
=\frac{40\,000}{3}\rub.
\]
В третьем году осталась треть основного долга:
\[
I_3=\frac{100\,000}{3}\cdot0{,}2
=\frac{20\,000}{3}\rub.
\]
Сумма процентов равна
\[
\Pi_2=20\,000+\frac{40\,000}{3}+\frac{20\,000}{3}
=40\,000\rub.
\]
Следовательно, вся выплата по второй схеме:
\[
T_2=100\,000+40\,000=140\,000\rub.
\]

\textbf{Сравнение.} Поскольку $140\,000<142\,417{,}58$, дешевле возвращать основной долг равными долями. Разница без промежуточного округления:
\[
\begin{aligned}
\Delta&=T_1-T_2\\
&\approx142\,417{,}58-140\,000\\
&=2\,417{,}58\rub.
\end{aligned}
\]

\subsection*{Проверка}
При второй схеме сумма трёх основных частей равна исходным $100\,000$ руб., а сумма процентов --- $40\,000$ руб.; значит, $T_2=140\,000$ руб. При первой схеме взнос одинаков каждый год и, согласно полученному уравнению, после третьего взноса остаток долга равен нулю. Сравниваются именно общие выплаты за одинаковый срок.

\Result{} меньшую общую выплату даёт погашение \textbf{равными частями основного долга}. Экономия составляет приблизительно \textbf{2\,417,58 руб.} при расчёте без промежуточного округления.

\subsection*{Задача на закрепление}
Кредит 99\,300 руб. взят на три года под 10\% годовых. Сравните погашение тремя равными ежегодными платежами и погашение равными частями основного долга. Найдите, какой способ даёт меньшую общую выплату и на сколько рублей.

\vfill
\begin{center}
\small Лёвин Артём Александрович эксклюзивно для Екатерины Гнедковой
\end{center}
\end{document}
